\section{Local metric frames and curvilinear machine coordinates}
\label{sec:moving-task-frames}

When saturation bends the flux map, the most meaningful directions may also
bend through current space. Let
\[
 \vect g_M=\nabla M,\qquad \vect g_P=\nabla P_{\mathrm{Cu}},\qquad
 \vect g_U=\nabla U^2.
\]
Each gradient is normal to its level surface. The Euclidean \(\vect g_M\)
maximizes torque change per Euclidean ampere step. If the scarce resource is
loss or voltage-to-current speed, a positive-definite metric \(\matr G\)
changes the answer to the natural-gradient direction
\begin{equation}
 \vect e_M=
 \frac{\matr G^{-1}\vect g_M}
 {\sqrt{\vect g_M^{\mathsf T}\matr G^{-1}\vect g_M}}.
 \label{eq:metric-torque-axis}
\end{equation}
With \(\matr G=\matr R\), it is the greatest torque change per incremental
loss length; with the dynamic metric of Part~III, it is the greatest torque
change per unit of local actuation difficulty.

\begin{figure}[tb]
\centering
\begin{tikzpicture}[>=Latex,font=\small,scale=1.05]
\draw[Purple,very thick] plot[smooth] coordinates {(-3,-1.1)(-1.7,-.2)(0,.1)(1.7,-.15)(3,-1.0)};
\node[Purple] at (2.7,-1.35){\(M=M_0\)};
\coordinate (P) at (0,.1);
\fill[BrickRed] (P) circle (2pt);
\draw[->,MidnightBlue,very thick] (P)--(0,2.0) node[above]{\(\nabla M\)};
\draw[->,ForestGreen!70!black,very thick] (P)--(1.1,1.45) node[right]{\(\matr G^{-1}\nabla M\)};
\draw[->,Orange,very thick] (P)--(2.0,-.18) node[right]{projected voltage direction};
\draw[dashed,gray] (-2.4,.6)--(2.5,.6);
\node[draw,rounded corners,fill=Orange!8,align=left] at (6.2,.55)
 {normal: change torque\\tangent 1: manage voltage\\tangent 2: remaining EESM allocation};
\end{tikzpicture}
\caption{A local task frame built from gradients. A physical metric changes
which direction is cheapest, so Euclidean orthogonality is not automatic.}
\label{fig:local-gradient-frame}
\end{figure}


\subsection{Metric Gram--Schmidt and a moving triad}

Metric orthogonality uses
\(\langle\vect x,\vect y\rangle_G=\vect x^{\mathsf T}\matr G\vect y\).
In two dimensions, remove from a trial vector \(\vect a\) its component along
a normalized \(\vect e_1\) by
\[
 \vect a_\perp=\vect a-\vect e_1(\vect e_1^{\mathsf T}\matr G\vect a).
\]
The same construction in the EESM gives a local triad
\[
 \matr E(\vect i)=[\vect e_M,\vect e_U,\vect e_N],
\]
where \(\vect e_U\) is the metric projection of
\(\matr G^{-1}\vect g_U\) into the constant-torque tangent plane and
\(\vect e_N\) completes the metric-orthonormal basis. Its meanings are
torque change, voltage management at constant torque, and remaining
allocation. The construction fails when \(\vect g_M=0\), when the projected
voltage gradient vanishes, or when gradients become nearly dependent.

For a saturated map
\(\boldsymbol\psi=\boldsymbol\phi(\vect i)\), the differential inductance
\(\matr L_{\mathrm{diff}}=\partial\boldsymbol\psi/\partial\vect i\) enters
the torque and voltage gradients. Consequently \(\matr E=\matr E(\vect i)\).
If local components are \(\vect q=\matr E^{-1}\dot{\vect i}\), then
\[
 \ddot{\vect i}=\matr E\dot{\vect q}
 +\underbrace{\dot{\matr E}\vect q}_{\text{motion of the basis}}.
\]
This is the three-dimensional analogue of the
\(-\dot\delta\matr J\vect x\) term in~\eqref{eq:arbitrary-frame-rate}.

\begin{figure}[tb]
\centering
\begin{tikzpicture}[>=Latex,font=\small,scale=.95]
\draw[->] (-3.3,0)--(3.4,0) node[right]{\(i_d\)};
\draw[->] (0,-2.5)--(0,2.7) node[above]{\(i_q\)};
\foreach \c in {.45,.85,1.25,1.65}
 {
  \draw[Purple!65,domain=.55:3,samples=70] plot ({\x},{\c/\x});
  \draw[Purple!35,domain=-3:-.55,samples=70] plot ({\x},{\c/\x});
 }
\foreach \a in {-1.0,-.5,0,.5,1.0}
 \draw[ForestGreen!65!black,domain=.45:2.05,samples=50]
   plot ({sqrt(2)*\x*exp(\a)/2},{sqrt(2)*\x*exp(-\a)/2});
\node[Purple] at (2.55,2.0){constant torque};
\node[ForestGreen!60!black] at (2.65,.62){allocation lines};
\node[align=left,draw,rounded corners,fill=MidnightBlue!6] at (6.0,.4)
 {state-dependent grid:\\axes bend with the flux map\\and acquire connection terms};
\end{tikzpicture}
\caption{A curvilinear grid that follows torque and allocation structure.
Under saturation the families deform and must be tabulated or evaluated from
the local flux Jacobian.}
\label{fig:curvilinear-grid}
\end{figure}


\subsection{Local arrows are not automatically global coordinates}

A smoothly varying triad is a useful differential frame, but its arrows need
not integrate to a globally consistent rectangular chart. In plain language:
walking first along the local torque arrow and then along the allocation arrow
may not reach the same point as taking the steps in reverse order. Vanishing
Lie brackets for the desired coordinate vector fields---or the corresponding
Frobenius involutivity condition for distributions---is the formal
integrability test. Torque itself is integrable by construction away from
critical points; an arbitrarily Gram--Schmidt-completed pair generally is not.
It should then be used as a local control basis, not advertised as a global
coordinate transformation.

Curvilinear charts are still useful. In the \(dq\) plane, torque curves can
label one family and a loss- or flux-based variable can label the crossing
family. The resulting grid follows the physics like polar coordinates follow
circles. Conformality is neither required nor sufficient for decoupled
dynamics. The square map from §cref{sec:conformal} preserves planar angles
away from its branch point but does not remove observer, input, or saturation
coupling; in three-dimensional EESM space, general metric-adapted
diffeomorphisms are more natural than forcing a conformal analogy.

Online identification could estimate
\(\partial M/\partial\vect i\) and
\(\partial\boldsymbol\psi/\partial\vect i\), then update the local frame
without a fixed parameter map. This is an exploratory adaptive extension, not
an established replacement for observers. Physical interpretability remains
a hard constraint: every learned axis must still state what motion along it
does, and continuity, excitation richness, sign alignment, and safety bounds
must be enforced.
